\paragraph{First-principles design.}
The reason to look at a neural operator here is not that the Sinkhorn path is conceptually broken. It is that the cost is wrong for repeated filtering. In \S5.3 the OT resampler costs $\mathcal{O}(LN^2)$ per resampling step, with $L$ Sinkhorn iterations and $N$ particles. A single filter run solves hundreds of nearby OT problems. The particle clouds change, but the structure of the problem does not. That is an amortisation opportunity.

I first considered the supervised route: store particle clouds, weights, and the Sinkhorn transport matrix, then train a map to predict the discrete transport. I did not take that route. The version here follows the Brenier / Monge--Amp\`ere route. It smooths the empirical source and target measures, writes the source as a uniform density and the target as a weighted density, and trains a conditional map $T(x;c)$ so that mass conservation is approximately satisfied. I used a direct-$T$ parameterisation rather than an input-convex potential $\phi$ because the loss needs $\det J_T$. If $T=\nabla\phi$, then differentiating the loss sees Hessians of the potential. With a direct map, $T(x;c)$ and $J_T(x;c)$ are produced together, so the determinant is a first-order object for the model.

\paragraph{What the Monge--Amp\`ere equation is here.}
The generic statement is simple. If a map $T:\mathbb{R}^d\to\mathbb{R}^d$ pushes a source density $p$ onto a target density $q$, then mass conservation gives the change-of-variables identity
\begin{equation}
    p(x) = q(T(x))\,\left|\det J_T(x)\right|.
\end{equation}
For squared Euclidean transport cost, Brenier's theorem says the optimal map is the gradient of a convex potential: $T=\nabla\phi$ with $\phi$ convex, under the usual absolute-continuity conditions. Substituting that structure gives the Monge--Amp\`ere form
\begin{equation}
    p(x) = q(\nabla\phi(x))\,\det \nabla^2\phi(x).
\end{equation}
That is the PDE route. I did not solve an explicit potential equation. I used the direct-map version of the same mass-conservation condition.

In our resampling problem, the inputs are particles $x_i\in\mathbb{R}^d$ with normalized importance weights $w_i$. I use the same scalar Gaussian bandwidth $h$ for the uniform and weighted clouds:
\begin{equation}
    K_h(u) = (2\pi h^2)^{-d/2}\exp\!\left(-\frac{\|u\|^2}{2h^2}\right).
\end{equation}
The source density is the uniform-smoothed cloud,
\begin{equation}
    p_h(x) = \frac{1}{N}\sum_{i=1}^{N} K_h(x-x_i),
\end{equation}
and the target density is the same cloud with the particle-filter weights,
\begin{equation}
    q_h(x) = \sum_{i=1}^{N} w_i K_h(x-x_i).
\end{equation}
The resampling-specific equation is therefore
\begin{equation}
    p_h(x) = q_h(T(x;c))\,\left|\det J_T(x;c)\right|.
\end{equation}
This is the condition imposed at collocation points sampled from the uniform density. The training loss evaluates $\log p_h$ at $x$, evaluates $\log q_h$ at $T(x;c)$, takes the determinant of $J_T(x;c)$, and squares the residual
\begin{equation}
    r(x;c) =
    \log p_h(x)
    - \log q_h(T(x;c))
    - \log\left|\det J_T(x;c)\right|,
    \qquad
    L_{\mathrm{MA}}(c) =
    \mathbb{E}_{x\sim p_h}\!\left[r(x;c)^2\right].
\end{equation}
The sign convention is immaterial because the residual is squared. What matters is that the map is trained against the same source density, target density, and Jacobian determinant that appear in the Monge--Amp\`ere equation.

\paragraph{Input and output of the operator.}
There are two inputs to the trunk map. The first is the point batch $x$ where the map is evaluated, with shape $(B,d)$ during collocation training or $(N,d)$ during resampling. The second is the context vector $c$ with shape $(C,)$, produced once from the weighted particle cloud. The set encoder does not see only coordinates and log-weights. It builds tokens
\begin{equation}
    t_i = [x_i,\log w_i,w_i] \in \mathbb{R}^{d+2},
\end{equation}
so the token tensor has shape $(N,d+2)$. A learned per-particle map is applied to those tokens. The outputs are pooled with $\sum_i w_i h_i$, then joined with weighted mean, upper-triangular covariance, $\log N$, $\log\mathrm{ESS}$, and $\log h$ before the final projection to $c$.

The trunk output is also explicit:
\begin{equation}
    T(x;c)\in\mathbb{R}^{B\times d},
    \qquad
    J_T(x;c)\in\mathbb{R}^{B\times d\times d}.
\end{equation}
The Jacobian is analytic. It is not recovered by differentiating each output coordinate during the loss. The loss uses $p_h(x)$, $q_h(T(x;c))$, and $\det J_T(x;c)$ from the same forward evaluation.

The module structure is the main architectural choice. Each conditional block keeps $W_m$ fixed with respect to context and lets $c$ shift biases and positive gates. The Jacobian contribution has the form $W_m^\top D_m W_m$, where $D_m$ is diagonal and nonnegative because the activation derivative and scale parameters are forced positive. This is why the Jacobian is symmetric PSD by construction. I chose that structure because the Monge--Amp\`ere loss asks for $\log\det J_T$ at many collocation points. Positive determinant could not be a soft afterthought.

\paragraph{Ansatz.}
The bet was threefold. First, a conditional map $T(x;c)$ should amortise across particle clouds if $c$ is produced by a permutation-invariant set encoder. Relabeling particles should not change the context except through the same set statistics, and the tests enforce that. Second, the Monge--Amp\`ere residual should be trainable without a Sinkhorn teacher. That is why the loss uses KDE density evaluations and a log-determinant term rather than MSE to a stored transport matrix. Third, the map should stay usable in a filter if the local geometry is controlled. The current regularisers are local-linearity, which checks that $T(x+\delta;c)$ is close to $T(x;c)+J_T(x;c)\delta$ on the KDE scale, and an eigenvalue lower-bound barrier, which penalizes $\lambda_{\min}(J_T)$ below a threshold.

\paragraph{Inference path.}
Once trained, the resampler does one deterministic forward pass. During a filter run I have a weighted cloud $\{(x_i,w_i)\}_{i=1}^N$. The filter-pipeline integration layer computes the same resolution-style bandwidth, encodes the whole cloud once to get $c$, and applies the same map to every particle:
\begin{equation}
    \tilde{x}_i = T(x_i;c), \qquad \tilde{w}_i = \frac{1}{N}.
\end{equation}
It also returns the local Jacobians $J_T(x_i;c)$ for covariance transport. There is no Sinkhorn loop at inference and no $N\times N$ transport matrix. That is the amortisation payoff if the map quality problem is solved.

\paragraph{Tests.}
The construction tests check the pieces that make the ansatz legal. One group checks the nonconditional map: output shapes, analytic-vs-automatic Jacobian agreement, symmetry, PSD structure, and near-identity residual initialization. A second group repeats the same checks with context fixed, then verifies context sensitivity and gradients back to $c$. A third group checks the set encoder: context shape, permutation invariance, variable $N$, differentiability, and response to distribution changes. The density tests check ESS, weighted moments, Silverman-style bandwidth sanity, KDE normalization, center-vs-far density ordering, uniform-density sampling, and differentiability through particles.

The loss-sanity tests ask whether the PDE residual is meaningful. One test uses the identity map with uniform weights, where the loss should be zero. Another uses a random map with uniform weights, where the loss should be nonzero. The oracle test builds the exact 1D monotone map by CDF matching; in the latest run the oracle loss was $1.19\times 10^{-9}$. That says the residual can be driven to zero by the correct map. If the learned model plateaus, it is not because the loss has an unavoidable KDE floor.

The quality tests are more severe. The interface smoke test checks shapes, uniform output weights, PSD local Jacobians, no transport matrix, no ancestor indices, and deterministic outputs after short training. The 1D quality test trains for $2000$ steps and asks whether the map beats identity on a minimal weighted cloud. The normalized-quality test repeats the same question with per-cloud normalization and fixed bandwidth. The eigenvalue-barrier tests check that the barrier is zero near identity, fires at the analytic hinge value for a mock small eigenvalue, and has finite gradients. The training and evaluation smoke tests check finite outputs, finite metrics, lower Monge--Amp\`ere residual after training, lower mean error on random-cloud evaluation data, and zero log-determinant failures.

\paragraph{What passed and failed.}
The suite collected $37$ tests: $35$ passed and $2$ failed. No tests were skipped or marked expected-failure. The normalized-quality run reduced loss from $10.74$ to $0.304$, and $h$ was stable at $0.35$, but late gradients spiked to $2852.64$ and held-out mean error got worse than identity ($2.25$ versus $1.00$). The stricter 1D resampler-quality run did not reach its final quality assertion because a checkpoint filename no longer matched the current neural-network library's naming rule. Its trace still showed huge initial loss, variable bandwidth, and large gradient spikes. This is not a benchmarked acceleration result. It is an implementation with useful construction tests and failing quality tests.

\paragraph{Open problem: initialization.}
The first failure is initialization. The residual map was supposed to start near identity, but random positive module gates made the residual large. The current workaround uses a near-one base scale and negative gate biases so the residual starts closer to identity. The remaining issue is input scaling: unnormalized clouds push bounded activations into saturation. The real fix is to make normalization part of the production training and inference contract, not just a diagnostic variant.

\paragraph{Open problem: quality.}
The practical blocker is quality. The infrastructure smoke test passes, but the map-quality tests do not yet prove that the operator moves particles toward the weighted target. The current workaround is to adjust synthetic training data, bandwidth policy, and normalization experiments. The real fix is still empirical: make the easy one-dimensional quality ladder pass before claiming anything about filters.

\paragraph{Open problem: gradient supervision.}
The Monge--Amp\`ere residual may be too weak by itself. It constrains the determinant and density ratio, not the full Jacobian. The proposed escalation is Gaussian-to-Gaussian pretraining with analytic maps and analytic Jacobians, or Sobolev-style Jacobian supervision. I would not add that until the quality and bandwidth failures are isolated, because it is a structural change to the training signal.

\paragraph{Open problem: Jacobian lower bound.}
The earlier map had a $J_T\ge I$ problem: a pure residual PSD map can expand but not compress. The current conditional map has the minimal workaround, a learned positive base scale $\lambda(c)$ so that $J_T\ge \lambda(c)I$ and $\lambda$ can fall below one. That fixes the most obvious impossibility, but it is still a global scale. If a cloud needs compression in the tails and expansion in the center, a single $\lambda(c)$ may be too crude.

\paragraph{Status.}
The filter-pipeline integration layer compiles and runs as a map-based resampler: it computes a resolution bandwidth, encodes the cloud, returns mapped particles with uniform weights, and exposes local Jacobians for covariance transport. It does not return an $N\times N$ transport matrix. I have not benchmarked it end-to-end against the Sinkhorn baseline, and the current tests do not justify any acceleration claim. The point of including this section is narrower: the design is clear, the failure modes are documented, and the test infrastructure is in place for a future implementer to pick it up without rediscovering the same traps.
